\documentclass[11pt,a4paper]{article}

% -----------------------------------------------------------------------------
% Shift Algebra: implementation-faithful mathematical tutorial
% Compile with: pdflatex MATHEMATICAL_TUTORIAL.tex (twice for the TOC)
% -----------------------------------------------------------------------------
\usepackage[utf8]{inputenc}
\usepackage[T1]{fontenc}
\usepackage[margin=1in]{geometry}
\usepackage{amsmath,amssymb,mathtools,bm}
\usepackage{booktabs,longtable,array,multirow}
\usepackage{xcolor}
\usepackage{graphicx}
\usepackage{enumitem}
\usepackage{listings}
\usepackage{hyperref}
\usepackage{microtype}

\hypersetup{
  colorlinks=true,
  linkcolor=blue!55!black,
  urlcolor=blue!60!black,
  citecolor=green!40!black,
  pdftitle={The Mathematics of Shift Algebra},
  pdfauthor={ShiftAlgebra repository}
}
\setlength{\parindent}{0pt}
\setlength{\parskip}{0.55em}
\setcounter{tocdepth}{3}
\allowdisplaybreaks

\newcommand{\R}{\mathbb{R}}
\newcommand{\E}{\mathbb{E}}
\newcommand{\ind}{\mathbf{1}}
\newcommand{\od}{\odot}
\newcommand{\dB}{\Delta B}
\newcommand{\dBh}{\widehat{\Delta B}}
\newcommand{\MB}{\texttt{M\_B}}
\newcommand{\MA}{\texttt{M\_A}}
\newcommand{\MC}{\texttt{M\_C}}
\newcommand{\LP}{L_p}
\newcommand{\softplus}{\operatorname{softplus}}
\newcommand{\sigmoid}{\operatorname{sigmoid}}
\newcommand{\sign}{\operatorname{sign}}
\newcommand{\diag}{\operatorname{diag}}
\newcommand{\rank}{\operatorname{rank}}
\newcommand{\BPE}{\operatorname{BPE}}
\newcommand{\ECE}{\operatorname{ECE}}
\newcommand{\pcacc}{\operatorname{pcacc}}
\newcommand{\acc}{\operatorname{acc}}
\newcommand{\conf}{\operatorname{conf}}
\newcommand{\margin}{\operatorname{margin}}
\newcommand{\drift}{\operatorname{drift}}

\lstset{
  basicstyle=\ttfamily\small,
  frame=single,
  breaklines=true,
  columns=fullflexible,
  backgroundcolor=\color{gray!5},
  keywordstyle=\color{blue!60!black}
}

\title{\textbf{The Mathematics of Shift Algebra}\\
\large An implementation-faithful tutorial from images to the $L_p$--Heisenberg operator product}
\author{Companion to the \texttt{ShiftAlgebra} repository}
\date{October 22, 2026}

\begin{document}
\maketitle

\begin{abstract}
This document develops the complete mathematical pipeline implemented in this
repository.  The task is to forecast how a fixed image classifier will behave
under an unseen ordered composition of distribution shifts, using only
measurements of the individual shifts and selected calibration pairs.  We begin
with images and the frozen classifier, derive the 518-dimensional shift
fingerprint and the five-dimensional behavior target, introduce the relevant
semigroup and non-commutative algebra, derive every accumulator in the headline
\MB{} model, state the exact neural readout and training objective, explain the
alternative methods and baselines, and finish with evaluation statistics,
worked examples, scope, and limitations.  Whenever the conceptual algebra and
the literal source code differ, both are stated explicitly.
\end{abstract}

\tableofcontents
\newpage

% =============================================================================
\section{The problem in one equation}
% =============================================================================

Let $f_\theta$ be a fixed, already-trained classifier.  A primitive shift $X$
is an image transformation such as brightness, blur, noise, contrast,
pixelation, or saturation.  A word
\[
S=(X_1,X_2,\ldots,X_m)
\]
means ``apply $X_1$ first, then $X_2$, and so on.''  The repository's goal is to
learn a forecasting map
\begin{equation}
\boxed{
\mathcal{F}_\Theta:\big((v_{X_1},\dB_{X_1}),\ldots,
(v_{X_m},\dB_{X_m})\big)
\longmapsto \dBh_S\in\R^5
}
\label{eq:problem}
\end{equation}
without running images transformed by the unseen composition $S$ through the
classifier at prediction time.

For each primitive $X_i$, the model receives:
\begin{enumerate}[label=(\roman*)]
  \item an unlabeled fingerprint $v_{X_i}\in\R^{518}$, measured from images
  transformed by that primitive alone; and
  \item a measured primitive behavior change $\dB_{X_i}\in\R^5$.
\end{enumerate}
The output is five numbers forecasting the changes in confidence margin,
representation drift, calibration error, per-class accuracy, and overall
accuracy under the full ordered word.

This is a \emph{zero-shot composition} claim, not a label-free claim.  Labels
are required once to measure each primitive behavior vector and selected
training-pair targets.  The zero-shot boundary says that the held-out composed
images and held-out composed labels are not inputs to the predictor.

\subsection{Why the problem is not ordinary regression}

Two structural facts make a naive regressor unreliable.

\textbf{Saturation.} If two shifts separately lower accuracy by $10$ and $15$
percentage points, their combination need not lower it by $25$ points.  The
same examples can be harmed twice, probabilities are bounded, and accuracy has
a floor.  Thus
\[
\dB_{X\to Y}\neq \dB_X+\dB_Y
\]
in general.

\textbf{Order.} A blur applied after noise can erase some of that noise, while
noise applied after blur survives.  Hence the physical transformations and
their effects can satisfy
\[
X\to Y\neq Y\to X,
\qquad
\dB_{X\to Y}\neq \dB_{Y\to X}.
\]
The method therefore needs a saturating symmetric component and a directional
antisymmetric component.

% =============================================================================
\section{Data-generating objects and notation}
% =============================================================================

Let the clean calibration sample be
\[
\mathcal{D}_{\mathrm{cal}}=\{(x_n,y_n)\}_{n=1}^{N},
\qquad x_n\in[0,1]^{C\times H\times W}.
\]
For the main CIFAR-10 V2 run, $C=3$, $H=W=32$, and $N=2000$ for calibration.
Held-out ground truth is measured separately on all $10{,}000$ CIFAR-10 test
images.

Write $\tau_{X,s}$ for primitive $X$ at severity $s$.  For a word
$S=((X_1,s_1),\ldots,(X_m,s_m))$, the implementation constructs
\begin{equation}
x_n^S=	au_{X_m,s_m}\!\left(\cdots
\tau_{X_2,s_2}(\tau_{X_1,s_1}(x_n))\cdots\right).
\label{eq:imageword}
\end{equation}
Every stage is clipped to $[0,1]$.  This convention matters: throughout this
document, the textual word $X_1\to X_2$ means $X_1$ is physically applied
first, whereas the corresponding linear operator product is $T_{X_2}T_{X_1}$.

Let the frozen classifier factor as
\[
f_\theta(x)=W h_\theta(x)+b,
\]
where $h_\theta(x)\in\R^{d_f}$ is the penultimate feature.  In the main
width-$0.25$ CIFAR ResNet-18, $d_f=128$.  The model input is normalized, but
pixel statistics below are computed on the unnormalized $[0,1]$ image.
Define
\[
f_n^S=h_\theta(\operatorname{Normalize}(x_n^S)),
\qquad
\ell_n^S=f_\theta(\operatorname{Normalize}(x_n^S)),
\qquad
p_n^S=\operatorname{softmax}(\ell_n^S).
\]

\subsection{The six image primitives}

The literal transformations in \path{src/corruptions_v2.py} are:
\begin{align}
\tau_{A,s}(x)&=\operatorname{clip}(s x), &&\text{brightness scaling},\\
\tau_{D,s}(x)&=\operatorname{clip}\big(\mu(x)+s(x-\mu(x))\big),
&&\text{contrast about the per-image mean},\\
\tau_{B,s}(x)&=k_s*_{\mathrm{sep}}x,
&&\text{separable Gaussian blur},\\
\tau_{C,s}(x)&=\operatorname{clip}(x+s\varepsilon),
&&\varepsilon\sim\mathcal{N}(0,I),\\
\tau_{P,s}(x)&=\operatorname{Up}_{\mathrm{nearest}}
\big(\operatorname{Down}_{\mathrm{bilinear}}(x;\operatorname{round}(H/s))\big),\\
\tau_{S,s}(x)&=\operatorname{clip}\big(g(x)+s(x-g(x))\big),
&&g(x)=\text{channel mean at each pixel}.
\end{align}
Thus, in the source code, $A$ is brightness and $D$ is contrast.  Gaussian
noise uses an order-independent deterministic condition seed.  Consequently,
$B\to C$ and $C\to B$ share the same underlying noise realization; their
difference measures order rather than a resampled noise field.

% =============================================================================
\section{From images to the 518-dimensional fingerprint}
% =============================================================================

This section answers exactly where the 518 inputs come from.

For a single primitive $X$, evaluate the clean and shifted versions of the
\emph{same} $N$ images.  Define paired feature displacements
\[
\delta f_n^X=f_n^X-f_n^0.
\]
For a vector-valued sample $a_1,\ldots,a_N$, let $\mu_n(a)$ and $\sigma_n(a)$
denote coordinatewise sample mean and standard deviation.  For the image
batch, let $\mu_{n,h,w}(x^X),\sigma_{n,h,w}(x^X)\in\R^3$ be the per-channel
statistics over images and spatial locations.  The fingerprint is
\begin{equation}
\boxed{
v_X=\begin{bmatrix}
\mu_n(f^X)\\
\sigma_n(f^X)\\
\mu_{n,h,w}(x^X)\\
\sigma_{n,h,w}(x^X)\\
\mu_n(\delta f^X)\\
\sigma_n(\delta f^X)
\end{bmatrix}
\in\R^{128+128+3+3+128+128}=\R^{518}.
}
\label{eq:518}
\end{equation}
No class label appears in Equation~\eqref{eq:518}.  Notice two details that are
easy to state incorrectly:
\begin{itemize}
  \item The penultimate feature is 128-dimensional in the main model, not 256.
  \item The raw displacement block is the \emph{signed mean}
  $\mu_n(\delta f^X)$.  A later six-moment summary takes its mean absolute
  magnitude, but the stored 128-vector itself is not replaced by an absolute
  value.
\end{itemize}

Why include both $f^X$ and $\delta f^X$?  The former records the absolute region
of feature space occupied by shifted data; the latter records movement from
the paired clean reference.  Two shifts can have similar final feature
variance but different displacement directions, so both are potentially
informative.

\subsection{The six scalar moments used as an encoder shortcut}

The encoder also compresses the six blocks into
\begin{equation}
q(v_X)=\begin{bmatrix}
\operatorname{mean}(v_{1:128}) &
\operatorname{mean}(v_{129:256}) &
\operatorname{mean}(v_{257:259}) &
\operatorname{mean}(v_{260:262}) &
\operatorname{mean}|v_{263:390}| &
\operatorname{mean}(v_{391:518})
\end{bmatrix}^{\!\top}\!\in\R^6.
\end{equation}
These correspond to global feature mean, feature spread, pixel mean, pixel
spread, absolute mean displacement magnitude, and displacement spread.

% =============================================================================
\section{The five-dimensional behavior target}
% =============================================================================

The full measured behavior vector has six coordinates; the first five are the
target and mean confidence is diagnostic only.

For each example, let $p_{n,(1)}^S\ge p_{n,(2)}^S$ be the largest two class
probabilities and $\hat y_n^S=\arg\max_c p_{n,c}^S$.  Then
\begin{align}
\margin(S)&=\frac1N\sum_{n=1}^N\left(p_{n,(1)}^S-p_{n,(2)}^S\right),\\
\drift(S)&=\left\|\frac1N\sum_{n=1}^N f_n^S-
\frac1N\sum_{n=1}^N f_n^0\right\|_2,\\
\pcacc(S)&=\frac1C\sum_{c=1}^{C}
\frac{\sum_n\ind[y_n=c]\ind[\hat y_n^S=c]}{\sum_n\ind[y_n=c]},\\
\acc(S)&=\frac1N\sum_n\ind[\hat y_n^S=y_n],\\
\conf(S)&=\frac1N\sum_n\max_c p_{n,c}^S.
\end{align}

Calibration uses ten confidence bins $I_r=[r/10,(r+1)/10)$:
\begin{equation}
\ECE(S)=\sum_{r=0}^{9}\frac{|I_r(S)|}{N}
\left|
\frac{1}{|I_r(S)|}\sum_{n\in I_r(S)}\ind[\hat y_n^S=y_n]
-
\frac{1}{|I_r(S)|}\sum_{n\in I_r(S)}\max_c p_{n,c}^S
\right|.
\end{equation}
Empty bins contribute zero.

Collect these as
\[
B(S)=\big[\margin(S),\drift(S),\ECE(S),\pcacc(S),\acc(S),\conf(S)\big].
\]
The clean baseline is $B_0=B(\varnothing)$, and the training/prediction target is
\begin{equation}
\boxed{
\dB_S=\big(B(S)-B_0\big)_{1:5}\in\R^5.
}
\label{eq:behavior}
\end{equation}
Typical signs are negative for margin and accuracy, positive for drift, and
mixed for ECE.  This is why the algebra uses a signed power transform rather
than assuming every coordinate is nonnegative.

% =============================================================================
\section{Algebraic foundations}
% =============================================================================

\subsection{Semigroups and words}

A semigroup is a set with a closed associative binary operation.  It need not
have an identity or inverse.  Image-corruption words naturally form a free
semigroup under concatenation:
\[
(XY)Z=X(YZ)=XYZ.
\]
The physical transformations may identify some words, but parenthesization of
sequential application is still irrelevant.  The modeling goal is an
approximately homomorphic representation
\[
\rho(X_1\cdots X_m)\approx \rho(X_m)\cdots\rho(X_1),
\]
where the reverse matrix order follows the ``first operation acts first''
convention for column vectors.

\subsection{Commutativity, commutators, and antisymmetry}

Associativity does not imply commutativity.  Matrices always satisfy
$(AB)C=A(BC)$ but usually satisfy $AB\ne BA$.  Their commutator is
\[
[A,B]=AB-BA,
\qquad [B,A]=-[A,B].
\]
A useful order channel should therefore be antisymmetric: swapping two shifts
should negate the channel while leaving symmetric severity information
unchanged.

\subsection{Why upper-triangular matrices}

Upper-triangular matrices can store running sums and pair interactions in a
closed product.  For example,
\[
T(w,c)=
\begin{pmatrix}1&0&w\\0&1&c\\0&0&1\end{pmatrix}
\quad\Longrightarrow\quad
T(w_2,c_2)T(w_1,c_1)=T(w_1+w_2,c_1+c_2).
\]
Likewise, a Heisenberg block
\begin{equation}
H(u,v)=
\begin{pmatrix}
1&u&\tfrac12uv\\
0&1&v\\
0&0&1
\end{pmatrix}
\label{eq:heis}
\end{equation}
stores a directional cross-term under multiplication.  All products in this
project are coordinatewise across parallel blocks; $uv$ means $u\od v$ when
$u$ and $v$ are vectors.

% =============================================================================
\section{The signed $L_p$ power law}
% =============================================================================

For behavior coordinate $k$, define a learned exponent
\begin{equation}
p_k=0.8+3.7\sigmoid(r_k),
\qquad p_k\in(0.8,4.5).
\label{eq:p}
\end{equation}
Ignoring numerical stabilization for one moment, define
\begin{equation}
\Phi_p(b)_k=\sign(b_k)|b_k|^{p_k},
\qquad
\Phi_p^{-1}(a)_k=\sign(a_k)|a_k|^{1/p_k}.
\end{equation}
The power-semigroup aggregate is
\begin{equation}
b_{L_p}(X_1,\ldots,X_m)=
\Phi_p^{-1}\!\left(\sum_{i=1}^m\Phi_p(\dB_{X_i})\right).
\label{eq:lp}
\end{equation}

For same-sign scalar effects $b_i$, Equation~\eqref{eq:lp} has magnitude
$(\sum_i|b_i|^p)^{1/p}$.  If $p>1$ and at least two $b_i$ are nonzero,
\begin{equation}
\max_i|b_i|\le \left(\sum_i|b_i|^p\right)^{1/p}
<\sum_i|b_i|.
\label{eq:normineq}
\end{equation}
For equal effects $|b_i|=\delta$, the $L_p$ result is $m^{1/p}\delta$ rather
than the additive $m\delta$.  This is the mathematical source of saturation.

A precise norm inequality worth remembering is
\begin{equation}
m^{1/p-1}\|b\|_1\le \|b\|_p\le \|b\|_1
\qquad (p\ge1).
\label{eq:correctnorm}
\end{equation}
The lower bound, not the upper bound, contains $m^{1/p-1}$.  Any statement with
$\|b\|_p\le m^{1/p-1}\|b\|_1$ is generally false (take $b=(1,0)$).  This
corrects the direction of the middle inequality printed in one prose proof in
\path{REVIEWER_PROOFING_V3.md}; the simpler strict sub-additivity conclusion in
Equation~\eqref{eq:normineq} remains valid.

\subsection{Associativity proof}

Define $a\oplus_p b=\Phi_p^{-1}(\Phi_p(a)+\Phi_p(b))$.  Because ordinary
addition is associative and $\Phi_p$ is bijective,
\begin{align}
(a\oplus_p b)\oplus_p c
&=\Phi_p^{-1}\big(\Phi_p(a)+\Phi_p(b)+\Phi_p(c)\big)\\
&=a\oplus_p(b\oplus_p c).
\end{align}
This is a coordinatewise semigroup isomorphic to $(\R^5,+)$.

\subsection{Literal numerical transform in the code}

With $\epsilon=10^{-6}$, the source uses
\begin{align}
\widetilde\Phi_p(b)&=\sign(b)
\left((|b|+\epsilon)^p-\epsilon^p\right),\\
\widetilde\Psi_p(a)&=\sign(a)(|a|+\epsilon)^{1/p}.
\end{align}
The matrix always sums $\widetilde\Phi_p(\dB_i)$ once and applies
$\widetilde\Psi_p$ once.  Thus the transformed accumulator is exactly
associative, but $\widetilde\Psi_p$ is not the exact algebraic inverse of
$\widetilde\Phi_p$.  Consequently, a recursively decoded binary operation has
a tiny $\epsilon$-level deviation from the ideal proof.  Also, if a learned
$p_k<1$, the coordinate is a signed quasi-norm and can be super-additive; the
empirical saturation interpretation applies to learned coordinates with
$p_k>1$.

% =============================================================================
\section{Shared encoder and decoder}
% =============================================================================

The shift encoder maps the fingerprint to a 64-dimensional latent vector:
\begin{align}
\widetilde z_X&=W_3\operatorname{ReLU}\!\left(
W_2\operatorname{ReLU}\!\left(
\operatorname{LN}(W_1v_X+b_1)\right)+b_2\right)+b_3,\\
z_X&=E(v_X)=\widetilde z_X+\tfrac12 W_q q(v_X)+\tfrac12b_q,
\qquad z_X\in\R^{64}.
\end{align}
The layer widths are $518\to256\to128\to64$.  The shared behavior decoder is
\begin{equation}
G(h)=W_5\operatorname{ReLU}(W_4h+b_4)+b_5,
\qquad 64\to128\to5.
\end{equation}
The encoder is behavior-anchored during training by requiring $G(E(v))$ to
predict observed behavior.  The 518-vector itself remains label-free; the
anchor loss is not.

% =============================================================================
\section{Headline method \MB: complete derivation}
% =============================================================================

\subsection{Dimensions and a crucial clarification}

The phrase ``a $3\times3$ matrix'' is shorthand.  The implementation uses many
parallel $3\times3$ upper-triangular blocks:
\[
d_b=5,
\qquad d_z=64,
\qquad r=8,
\qquad K_{\mathrm{ord}}=d_b+2r=21.
\]
Define
\[
d_w=d_b+2d_z+2K_{\mathrm{ord}}=175,
\qquad
L=d_w+d_z+K_{\mathrm{ord}}=260.
\]
Therefore $T_X$ is stored as an array in $\R^{260\times3\times3}$.  Equivalently,
one could view it as a block-diagonal $780\times780$ matrix, but the code never
needs to materialize that dense matrix.  There are 175 translation blocks, 64
symmetric Heisenberg blocks, and 21 directional Heisenberg blocks.

\subsection{Potential and amplitude}

Build the operator input
\[
o_X=[z_X;\dB_X]\in\R^{69}.
\]
Two MLPs produce a potential $\kappa_X\in\R^{21}$ and raw amplitude
$a_X\in\R^{21}$.  The positive amplitude $\pi_X$ is
\begin{equation}
\pi_{X,k}=
\begin{cases}
(|\dB_{X,k}|+0.05)\softplus(a_{X,k}+0.55),&1\le k\le5,\\
\softplus(a_{X,k}),&6\le k\le21.
\end{cases}
\label{eq:pi}
\end{equation}
Set
\begin{equation}
u_X=\kappa_X\od\pi_X,
\qquad q_X=\pi_X.
\label{eq:uq}
\end{equation}
We use $q_X$ rather than $v_X$ for the second Heisenberg generator here to
avoid confusing it with the 518-dimensional fingerprint.

The directional pair term then has a potential-difference form:
\begin{equation}
\frac12(u_Y\od q_X-u_X\od q_Y)
=\frac12\pi_X\od\pi_Y\od(\kappa_Y-\kappa_X).
\label{eq:potentialdiff}
\end{equation}
This imposes a transitive structure: the sign comes from the ordering of
learned scalar-like susceptibilities rather than an unrelated parameter for
every pair.

\subsection{Primitive accumulators}

For one primitive, concatenate
\begin{equation}
w_X=\left[
\widetilde\Phi_p(\dB_X);\ z_X;\ z_X^{\od2};\ u_X;\ u_X\od q_X
\right]\in\R^{175}.
\label{eq:w}
\end{equation}
For each coordinate $w_{X,\ell}$, create a translation block
\begin{equation}
A_\ell(X)=
\begin{pmatrix}
1&0&w_{X,\ell}\\
0&1&1\\
0&0&1
\end{pmatrix}.
\end{equation}
For each latent coordinate $j$, create a symmetric block
\begin{equation}
H_j^{\mathrm{sym}}(X)=
\begin{pmatrix}
1&z_{X,j}&\tfrac12 z_{X,j}^2\\
0&1&z_{X,j}\\
0&0&1
\end{pmatrix}.
\end{equation}
For each directional coordinate $k$, create
\begin{equation}
H_k^{\mathrm{ord}}(X)=
\begin{pmatrix}
1&u_{X,k}&\tfrac12u_{X,k}q_{X,k}\\
0&1&q_{X,k}\\
0&0&1
\end{pmatrix}.
\end{equation}
The complete $T_X$ is the direct sum of all these blocks.

\subsection{Ordered product}

For $S=(X_1,\ldots,X_m)$ in application order,
\begin{equation}
\boxed{T_S=T_{X_m}T_{X_{m-1}}\cdots T_{X_1}.}
\label{eq:orderedprod}
\end{equation}
Every stored $3\times3$ block is multiplied only by its corresponding block.
Apply each block to $s_0=(0,0,1)^\top$, which is just taking its third column.
The translation blocks return
\begin{equation}
\sum_iw_{X_i}
=\left[
\sum_i\widetilde\Phi_p(\dB_{X_i});\
\sum_i z_{X_i};\
\sum_i z_{X_i}^{\od2};\
\sum_i u_{X_i};\
\sum_i u_{X_i}\od q_{X_i}
\right],
\end{equation}
and their second state coordinate returns the count $m$.

\subsection{Deriving the Heisenberg top-right entry}

For $H_i=H(u_i,q_i)$, direct multiplication gives
\[
H_2H_1=
\begin{pmatrix}
1&u_1+u_2&
\tfrac12u_1q_1+\tfrac12u_2q_2+u_2q_1\\
0&1&q_1+q_2\\
0&0&1
\end{pmatrix}.
\]
Induction yields
\begin{equation}
(H_m\cdots H_1)_{0,2}
=\frac12\sum_{i=1}^m u_iq_i+
\sum_{1\le i<j\le m}u_jq_i.
\label{eq:heistop}
\end{equation}
Let $N_2=\binom{m}{2}$, with the code clamping it to at least one.  This one
identity produces both interaction channels.

\subsection{Symmetric all-pairs interaction}

Set $u_i=q_i=z_i$ in the symmetric blocks.  Subtract the stored self-square
terms:
\begin{align}
\Gamma_z(S)
&=\frac{(H_S^{\mathrm{sym}})_{0,2}
-\tfrac12\sum_i z_i^{\od2}}{N_2}\\
&=\frac1{N_2}\sum_{i<j}z_i\od z_j.
\label{eq:gamma}
\end{align}
This is invariant under every permutation of the input shifts.

\subsection{Antisymmetric order interaction}

Expand half the product of sums:
\begin{equation}
\frac12\left(\sum_i u_i\right)\od\left(\sum_iq_i\right)
=\frac12\sum_i u_i\od q_i
+\frac12\sum_{i<j}(u_i\od q_j+u_j\od q_i).
\end{equation}
Subtracting this from Equation~\eqref{eq:heistop} gives
\begin{align}
\Omega_{\mathrm{ord}}(S)
&=\frac{(H_S^{\mathrm{ord}})_{0,2}
-\tfrac12(\sum_i u_i)\od(\sum_iq_i)}{N_2}\\
&=\frac1{N_2}\sum_{i<j}\frac12
(u_j\od q_i-u_i\od q_j)\\
&=\frac1{N_2}\sum_{i<j}\frac12
\pi_i\od\pi_j\od(\kappa_j-\kappa_i).
\label{eq:omega}
\end{align}
For the fully reversed word $S^{\mathrm{rev}}$,
\begin{equation}
\Gamma_z(S^{\mathrm{rev}})=\Gamma_z(S),
\qquad
\Omega_{\mathrm{ord}}(S^{\mathrm{rev}})=-\Omega_{\mathrm{ord}}(S).
\end{equation}
This is an exact identity, not a learned preference.

\subsection{Exact readout used by the source}

The matrix state is reduced to
\begin{align}
b_{L_p}&=\widetilde\Psi_p\!\left(\sum_i
\widetilde\Phi_p(\dB_{X_i})\right),\\
\bar z&=\frac1m\sum_i z_{X_i},\\
\Gamma_z&\text{ from Equation~\eqref{eq:gamma}},\\
\Omega&\text{ from Equation~\eqref{eq:omega}}.
\end{align}
The current V3 source intentionally sets the former linear pair-sum input to
zero:
\[
z_{\mathrm{pair\_scale}}=0.
\]
This prevents a network trained on pairs from relying on a first-order sum that
scales incorrectly on triples.  Define
\begin{equation}
s_{\mathrm{sym}}=[0;\Gamma_z;b_{L_p}].
\end{equation}
Then
\begin{align}
h_{\mathrm{sym}}&=\bar z+S_{\mathrm{MLP}}(s_{\mathrm{sym}}),\\
g&=0.5+\sigmoid\big(G_{\mathrm{gate}}([\bar z;b_{L_p}])\big)
\in(0.5,1.5)^5,\\
b_{\mathrm{beh}}&=g\od b_{L_p}+R_{\mathrm{sym}}(s_{\mathrm{sym}}),\\
b_{\mathrm{sym}}&=\tfrac12b_{\mathrm{beh}}+	frac12G(h_{\mathrm{sym}}),\\
a_{\mathrm{ord}}&=0.25\,\Omega_{1:5}+W_{\mathrm{ord}}^{(b)}\Omega,\\
\boxed{\dBh_S&=b_{\mathrm{sym}}+a_{\mathrm{ord}}.}
\label{eq:mbreadout}
\end{align}
All MLP residual output layers are initialized at zero, so optimization starts
near the interpretable $L_p$ base instead of an arbitrary black-box map.  The
source class declares a latent map named \texttt{W\_ord}; the current \MB{}
\texttt{\_readout} does not use that map.  Order enters the behavior output
through the direct first-five-coordinate term and \texttt{ord\_beh}, shown as
$W_{\mathrm{ord}}^{(b)}$ above.

\subsection{What is exactly associative}

Matrix multiplication gives
\begin{equation}
T_{X_3}(T_{X_2}T_{X_1})=(T_{X_3}T_{X_2})T_{X_1}
\label{eq:assoc}
\end{equation}
up to floating-point roundoff.  Therefore any single readout applied to the
same total product is parenthesization-independent.  This is what the
repository's \texttt{associativity\_error} diagnostic verifies for \MB{}.
It does not claim that every arbitrary learned binary decoder is associative.

\subsection{The exact three-shift Jordan-style readout}

There is an additional implementation detail for $m=3$.  Let
$\mathcal{P}((h_1,b_1),(h_2,b_2))$ denote \texttt{pair\_forward}, and let
$(h_\star,b_\star)$ be the single-pass readout of $T_3T_2T_1$.  Compute both
binary trees
\begin{align}
(h_{12},b_{12})&=\mathcal{P}((z_1,\dB_1),(z_2,\dB_2)),\\
(h_L,\_)&=\mathcal{P}((h_{12},b_{12}),(z_3,\dB_3)),\\
(h_{23},b_{23})&=\mathcal{P}((z_2,\dB_2),(z_3,\dB_3)),\\
(h_R,\_)&=\mathcal{P}((z_1,\dB_1),(h_{23},b_{23})).
\end{align}
The source then returns
\begin{align}
h&=0.12h_\star+0.44(h_L+h_R),\\
\dBh&=0.12b_\star+0.44\big(G(h_L)+G(h_R)\big)+0.44a_{\mathrm{ord},\star}.
\label{eq:jordan}
\end{align}
Averaging both parenthesizations removes a preference for one binary tree while
retaining the ordered matrix commutator.  This extra readout is not itself a
plain matrix product.  The strongest precise statement is: the operator core
in Equation~\eqref{eq:assoc} is exactly associative, and the triple predictor
explicitly symmetrizes its two recursive decoding trees.  Reported
\texttt{assoc\_err} measures the former.

\subsection{Out-of-support safeguard for a new primitive}

At inference, the code compares each primitive fingerprint with the primitive
support represented in training.  Let $d$ be the maximum relative nearest
support distance.  If $d>0.08$, it predicts both forward and reversed words,
\[
b_{\mathrm{sym}}=\tfrac12(\dBh_S+\dBh_{S^{\mathrm{rev}}}),
\qquad
b_{\mathrm{ord}}=\tfrac12(\dBh_S-\dBh_{S^{\mathrm{rev}}}),
\]
and shrinks toward the Stage-0 $L_p$ law:
\begin{equation}
\dBh_S=b_{L_p}+g_d(b_{\mathrm{sym}}-b_{L_p})+c_d b_{\mathrm{ord}},
\end{equation}
where
\[
g_d=\exp\{-150\max(0,d-0.08)^2\},
\qquad
c_d=\begin{cases}1,&d>0.25,\\0.15,&0.08<d\le0.25.\end{cases}
\]
This guard is especially relevant to primitive-holdout and counterfactual-shift
tests.  It is part of the implemented predictor, not part of the abstract
Heisenberg theorem.

% =============================================================================
\section{Alternative proposed methods}
% =============================================================================

\subsection{\MA: behavior-space anchored algebra}

\MA{} uses the same $L_p$ base but computes interactions explicitly rather
than extracting them from a matrix product.  For a word of length $m$,
\begin{align}
\bar z&=\frac1m\sum_i z_i,\\
\Gamma_z&=\frac1{\binom m2}\sum_{i<j}z_i\od z_j,\\
\Omega&=\frac1{\binom m2}\sum_{i<j}
\frac12\pi_i\od\pi_j\od(\kappa_j-\kappa_i),\\
z_{\mathrm{scale}}&=\frac2m\sum_i z_i.
\end{align}
Its symmetric input is $[z_{\mathrm{scale}};\Gamma_z;b_{L_p}]$.  It predicts
both a latent representation and behavior:
\begin{align}
h&=\bar z+H_{\mathrm{MLP}}([z_{\mathrm{scale}};\Gamma_z;b_{L_p}])
+W_h\Omega,\\
\widetilde b&=g(\bar z,b_{L_p})\od b_{L_p}
+R([z_{\mathrm{scale}};\Gamma_z;b_{L_p}])+W_b\Omega,\\
\dBh&=\tfrac12\widetilde b+\tfrac12G(h).
\end{align}
For triples, its sequence function also averages the direct all-pairs result
with both binary trees:
\[
h\leftarrow0.4h+0.3(h_L+h_R),
\qquad
\dBh\leftarrow0.4\dBh+0.3(b_L+b_R).
\]
\MA{} is therefore ``Lie-inspired'' through an antisymmetric potential
difference.  The literal implementation is not a truncated
Baker--Campbell--Hausdorff series.

\subsection{\MC: ordered cross-attention fold}

For a pair, form tokens
\[
t_i=[z_i;\dB_i;e_i],
\]
where $e_i\in\R^8$ is a learned position embedding.  Compute
\begin{align}
Q_i&=W_Qt_i,\quad K_i=W_Kt_i,\quad V_i=W_Vt_i,\\
A_{ij}&=\frac{\exp(Q_i^\top K_j/\sqrt{32})}
{\sum_{j'}\exp(Q_i^\top K_{j'}/\sqrt{32})},\\
o_i&=\sum_jA_{ij}V_j.
\end{align}
The ordered feature is
\[
r=[o_1;o_2;\dB_1+\dB_2;\dB_1-\dB_2].
\]
A network produces separate gates $g_1,g_2\in(0,1.5)^5$ and residual $r_b$:
\begin{equation}
\dBh=g_1\od\dB_1+g_2\od\dB_2+r_b,
\qquad
h=\tfrac12(z_1+z_2)+r_h.
\end{equation}
Longer words are left-folded.  Removing position tokens yields a deliberately
permutation-invariant ablation; its order sign accuracy collapses because it
cannot distinguish $X\to Y$ from $Y\to X$.

\subsection{Set-attention variant}

The \texttt{M\_A\_att} method uses a learned query to pool a whole token set.
With keys $k_i$, values $r_i$, and query $q$,
\[
\alpha_i=\frac{\exp(k_i^\top q/\sqrt{64})}
{\sum_j\exp(k_j^\top q/\sqrt{64})},
\qquad c=\sum_i\alpha_i r_i.
\]
The context $c$ controls per-shift retention gates.  It is flexible and
sub-additive but lacks the exact matrix-product decomposition of \MB{}.

% =============================================================================
\section{Baselines and decisive ablations}
% =============================================================================

Understanding the controls is necessary to understand what the experiment
actually tests.

\begin{longtable}{>{\raggedright\arraybackslash}p{0.20\textwidth}
                  >{\raggedright\arraybackslash}p{0.72\textwidth}}
\toprule
Method & Mathematical rule or purpose\\
\midrule
\endfirsthead
\toprule
Method & Mathematical rule or purpose\\
\midrule
\endhead
\texttt{M0\_meanpool} & $h=m^{-1}\sum_i z_i$, $\dBh=G(h)$.  Commutative and unable to encode order.\\
\texttt{M1\_behadd} & $\dBh=\sum_i\dB_i$.  This is the $p=1$ behavior baseline and commonly overshoots triples.\\
\texttt{M1\_emb} & $h=\sum_i z_i$, $\dBh=G(h)$.\\
\texttt{M\_Lp\_only} & Equation~\eqref{eq:lp} with only five learned exponents.  It isolates saturation and has no order channel.\\
\texttt{M\_PowerOffset} & $\alpha\od b_{L_p}+\beta$ with per-coordinate $p$, scale, and offset: 15 scalar parameters.\\
\texttt{M\_Max\_pool} & $\sign(\sum_i\dB_i)\od\max_i|\dB_i|$, the $p\to\infty$ envelope.\\
\texttt{M\_Ord\_noLp} & Full \MB{} with $p\equiv1$.  It tests whether the order network alone can replace saturation.\\
\texttt{M\_B\_nocomm} & Full symmetric \MB{} readout with $\Omega_{\mathrm{ord}}$ zeroed.  It isolates the commutator.\\
\texttt{M\_LowRankBilinear} & $h=0.35(z_A+z_B)+0.65(Az_A+Bz_B)+W(Uz_A\od Vz_B)$, left-folded.\\
\texttt{M\_MLP\_matched} & Ordered black-box MLP on $[z_A;\dB_A;z_B;\dB_B]$, parameter-matched as closely as practical and left-folded.\\
\texttt{M5\_v1} & $S([z_A+z_B;z_A\od z_B])+W(z_A\od Uz_B-z_B\od Uz_A)$; the earlier unconstrained latent algebra.\\
\texttt{MEMO\_control} & Exact lookup on seen conditions and mean prediction otherwise.  It tests memorization.\\
\texttt{PERM\_control} & Trains with permuted pair targets.  It tests whether results survive destroyed input--target alignment.\\
\bottomrule
\end{longtable}

Two V3 ablations are especially diagnostic:
\[
\BPE(\MB)=0.0152,
\quad
\BPE(\text{no }L_p)=0.1170,
\quad
\BPE(\text{no commutator})=0.0364.
\]
Removing the commutator increases order-pair BPE from $0.0297$ to $0.0912$ and
reduces order sign match from $0.90$ to $0.16$.  These comparisons isolate the
two intended mathematical roles rather than merely comparing two differently
sized neural networks.

% =============================================================================
\section{Training mathematics}
% =============================================================================

\subsection{Observed and held-out conditions}

There are six primitives and five severity levels, giving 30 primitive
conditions.  The main Regime-II training set has 25 ordered pair series, each
at five matched severity levels, hence
\[
30+25\cdot5=155
\]
observed conditions.  The pair graph contains:
\begin{itemize}
  \item four one-direction order-test chords $BC,BA,CS,PC$;
  \item their reverses $CB,AB,SC,CP$ held out;
  \item ten bidirectional training families $BP,AC,AD,AP,DS,AS,BD,BS,DP,PS$;
  \item one additional one-direction family $CD$.
\end{itemize}
The held-out set has 100 triples ($\binom63\cdot5$), 20 reverse-order pairs,
and 3 off-grid conditions, for 123 total.  Regime III adds six triple families
to training.

\subsection{Stage-0 exponent fit}

Before gradient training, each $p_k$ is initialized by grid search on observed
pairs only, over 143 values from $0.85$ to $4.40$.  The grid-search predictor is
\begin{equation}
\widehat b_k(p)=\sign(b_{1k}+b_{2k})
\left(|b_{1k}|^p+|b_{2k}|^p\right)^{1/p}.
\end{equation}
No triple target enters this selection.  The selected value is transformed
through Equation~\eqref{eq:p} to initialize the unconstrained parameter $r_k$.
The $r_k$ parameters are excluded from weight decay, preventing regularization
from pulling all $p_k$ back toward the sigmoid midpoint.

\subsection{Weighted behavior loss}

For prediction $\hat y$ and target $y$, the source defines
\begin{equation}
\mathcal{L}_w(\hat y,y)=
\frac1N\sum_{n=1}^N\sum_{k=1}^{5}w_k(\hat y_{nk}-y_{nk})^2,
\qquad
w=(1,1,1,1,0.3).
\label{eq:wmse}
\end{equation}
Accuracy is downweighted because it floors more readily than margin, drift,
ECE, or per-class accuracy.

\subsection{Main objective}

Let $z_i=E(v_i)$ for primitives, let $z_{ij}^{\mathrm{direct}}=E(v_{ij})$ be the
encoding of an \emph{observed training pair}, and let
$(h_{ij},\hat y_{ij})$ be the composition of the two primitive inputs.  The
main terms are
\begin{align}
\mathcal{L}_{\mathrm{pair}}&=\mathcal{L}_w(\hat y_{ij},y_{ij}),\\
\mathcal{L}_{\mathrm{anchor}}&=\tfrac12\left[
\mathcal{L}_w(G(z_i),y_i)+
\mathcal{L}_w(G(z_{ij}^{\mathrm{direct}}),y_{ij})\right],\\
\mathcal{L}_{\mathrm{repr}}&=\tfrac12\|h_{ij}-z_{ij}^{\mathrm{direct}}\|_2^2.
\end{align}
For bidirectionally observed pairs, define
\[
\Delta\hat y_{ij}=\hat y_{ji}-\hat y_{ij},
\qquad
\Delta y_{ij}=y_{ji}-y_{ij}.
\]
The additional order terms are
\begin{equation}
\mathcal{L}_{\mathrm{order}}
=2\mathcal{L}_w(\Delta\hat y,\Delta y)
+0.25\mathcal{L}_w(d_{\mathrm{sym}},0)
+0.10\mathcal{L}_w(a_{\mathrm{ord}},0).
\end{equation}
The latter two terms partition variance between symmetric and antisymmetric
channels while the difference loss gives the order channel a supervised sign.
For adjacent severities of the same primitive, \MB{} also penalizes a false
self-commutator:
\begin{equation}
\mathcal{L}_{\mathrm{semi}}=0.05\,\E\|\Omega(X_{s_a},X_{s_b})\|_2^2.
\end{equation}
In Regime III, a direct triple behavior loss is added.  Schematically,
\begin{equation}
\mathcal{L}=\mathcal{L}_{\mathrm{pair}}+
\mathcal{L}_{\mathrm{anchor}}+\mathcal{L}_{\mathrm{repr}}+
\mathcal{L}_{\mathrm{order}}+
\mathcal{L}_{\mathrm{semi}}+
\mathcal{L}_{\mathrm{triple}}.
\label{eq:loss}
\end{equation}
Only applicable terms are present for a given method/regime.

An implementation detail: \path{src/config_v2.py} defines
\texttt{LAMBDA\_REPR=0.2}, but the current training loop uses that value as an
enable/disable check and applies an effective coefficient $0.5$ to the
representation MSE.  Equation~\eqref{eq:loss} reflects the executable source.
Optimization uses Adam with learning rate $10^{-3}$, 500 steps by default, and
weight decay $10^{-4}$ on all regular parameters but zero on \texttt{raw\_p}.

\subsection{What pair data is and is not used at test time}

During training, an observed pair's direct fingerprint $v_{ij}$ provides a
representation-alignment target and its measured behavior $y_{ij}$ provides
supervision.  During zero-shot evaluation of a held-out composition, prediction
uses only the primitive cache
\[
\{(v_{X_i},\dB_{X_i})\}_{i=1}^m.
\]
The held-out composed fingerprint and target are loaded only to score the
prediction or compute the explicitly labeled non-zero-shot oracle
$G(E(v_{\mathrm{held}}))$.

% =============================================================================
\section{Controlled synthetic mathematics}
% =============================================================================

The synthetic suite checks mechanisms under known ground truth before the
vision benchmark.

\subsection{S1: additive and commuting}

The 64-dimensional feature space is divided into disjoint coordinate blocks.
Each primitive changes only its own block.  Operators on disjoint blocks
commute, so $T_XT_Y=T_YT_X$, and their behavior effects are nearly additive.
This checks that interaction-aware models do not invent large commutators when
none exist.

\subsection{S2: interacting and saturating}

A state has displacement $d$ and attenuation $a$.  Each shift updates
schematically as
\begin{align}
d&\leftarrow d+a\od s\delta_X,\\
a&\leftarrow a\od\exp(-0.55s\sigma_X),\\
x'&=a\od x+1.4\tanh(d/1.4).
\end{align}
The exponential attenuation and hyperbolic tangent create diminishing returns
and pair interactions.  This isolates why an $L_p$-like base can beat linear
addition.

\subsection{S3: non-commutative filter then perturb}

Ignoring its final saturation for clarity, one step is affine:
\begin{equation}
x\longmapsto (I-sD_X)x+s\delta_X.
\end{equation}
For two shifts,
\begin{align}
T_Y(T_X(x))&=(I-s_YD_Y)(I-s_XD_X)x
 +(I-s_YD_Y)s_X\delta_X+s_Y\delta_Y,\\
T_X(T_Y(x))&=(I-s_XD_X)(I-s_YD_Y)x
 +(I-s_XD_X)s_Y\delta_Y+s_X\delta_X.
\end{align}
Even if the linear damping matrices happened to commute, the affine terms need
not: the second filter damps the first perturbation, but not its own newly
written perturbation.  This is the synthetic analogue of blur/noise order.

% =============================================================================
\section{Evaluation metrics and statistical inference}
% =============================================================================

Let $P_{s,h,k}$ be the prediction from random seed $s$ on held-out condition
$h$ and behavior coordinate $k$, and let $Y_{h,k}$ be ground truth.  The
primary behavior prediction error is
\begin{equation}
\BPE=\frac1{S H\cdot4}\sum_{s=1}^{S}\sum_{h=1}^{H}
\sum_{k=1}^{4}(P_{s,h,k}-Y_{h,k})^2.
\label{eq:bpe}
\end{equation}
Accuracy is reported separately as
\[
\operatorname{AccMAE}=\frac1{SH}\sum_{s,h}|P_{s,h,5}-Y_{h,5}|.
\]
The first four coordinates are primary because accuracy can floor under strong
corruption.

\subsection{Ranking and order metrics}

Spearman correlation is computed per seed between predicted and measured
margin (and separately accuracy), then averaged.  For each reverse-order pair,
order sign match is tested on drift:
\begin{equation}
\sign(P_{Y\to X,2}-P_{X\to Y,2})
\stackrel{?}{=}
\sign(Y_{Y\to X,2}-Y_{X\to Y,2}).
\end{equation}
Commutator correlation is the Spearman correlation between predicted and true
magnitudes
\[
\|\dB_{Y\to X}-\dB_{X\to Y}\|_{2,1:4}
\]
across pair families and severities.  Order error is the mean absolute error
between those predicted and true magnitudes.

\subsection{Interaction residual error}

For a composition $S$, define its non-additive residual
\begin{equation}
\gamma(S)=\dB_S-\sum_{X_i\in S}\dB_{X_i}.
\end{equation}
Interaction residual error (IRE) is the MSE between predicted and true
$\gamma(S)$ over the first four coordinates, reported separately for pairs and
triples.  It tests whether a model captures the interaction rather than merely
ranking overall severity.

\subsection{Hierarchical statistical tests}

The V3 analysis avoids treating every seed--severity point as fully
independent:
\begin{enumerate}
  \item It averages condition error within each of five seeds and reports a
  paired seed-level $t$ test and confidence interval.
  \item It averages seeds, groups 123 conditions into 25 composition families
  (20 triple families, 4 reverse-pair families, and 1 off-grid family), then
  cluster-bootstraps families.
  \item It computes condition-level paired Wilcoxon signed-rank tests and
  applies Holm--Bonferroni correction across 11 primary comparisons.
\end{enumerate}
For a family of $M$ raw $p$-values sorted as
$p_{(1)}\le\cdots\le p_{(M)}$, the Holm adjusted values are the monotone
step-down quantities
\begin{equation}
\widetilde p_{(i)}=
\max_{j\le i}\min\{1,(M-j+1)p_{(j)}\}.
\end{equation}

% =============================================================================
\section{Worked examples}
% =============================================================================

\subsection{A numerical $L_p$ example}

Suppose two primitive accuracy deltas are $b_X=-0.10$ and $b_Y=-0.15$ and
$p=2$.  Linear addition predicts $-0.25$.  The power law predicts
\[
b_{L_2}=-\sqrt{0.10^2+0.15^2}=-0.1803.
\]
It is larger in magnitude than either single drop but smaller than their sum,
which is exactly the desired diminishing-return geometry.

For three equal $-0.10$ shifts,
\[
b_{L_2}=-\sqrt{3}\,(0.10)=-0.1732,
\]
whereas addition predicts $-0.30$ and averaging predicts $-0.10$.

\subsection{A numerical Heisenberg example}

Take one directional coordinate with
\[
\kappa_X=0.2,\quad\pi_X=0.5,
\qquad
\kappa_Y=1.0,\quad\pi_Y=0.4.
\]
Then $u_X=0.1,q_X=0.5,u_Y=0.4,q_Y=0.4$, and
\begin{align}
\Omega(X\to Y)
&=\tfrac12(u_Yq_X-u_Xq_Y)\\
&=\tfrac12(0.4\cdot0.5-0.1\cdot0.4)=0.08.
\end{align}
Swapping the shifts gives $-0.08$ exactly.  No symmetric interaction has to
change when the order changes.

\subsection{Why the product remains associative but order-sensitive}

For three primitive matrices,
\[
(T_ZT_Y)T_X=T_Z(T_YT_X),
\]
so changing parentheses does nothing.  But generally
\[
T_ZT_YT_X\ne T_XT_YT_Z,
\]
so reversing the physical order can change the result.  Associativity and
commutativity are independent properties; \MB{} needs the first and rejects
the second.

% =============================================================================
\section{What is universal and what is dataset-dependent}
% =============================================================================

The \emph{form} of the law is reusable:
\begin{itemize}
  \item signed $L_p$ accumulation for saturation;
  \item parallel upper-triangular Heisenberg blocks for all-pairs interactions;
  \item ordered matrix multiplication for non-commutativity and associativity;
  \item a symmetric readout plus an antisymmetric order readout.
\end{itemize}
The same equations were run on CIFAR-10 with a ResNet, CIFAR-10 with a
ViT/PatchMixer, Fashion-MNIST, and CIFAR-100 without changing the algebraic
form.

The \emph{parameters} are deployment-dependent.  A new dataset/classifier pair
has a different feature coordinate system and behavior scale, so it requires
new primitive fingerprints, primitive behavior measurements, and calibration
of $E,G,p,\kappa,\pi$, and the readout networks.  Thus \MB{} is not one frozen
universal matrix for every dataset.  It is a reusable structural model that is
fit per dataset--classifier deployment.

After such calibration, the project tests two levels of generalization:
\begin{enumerate}
  \item new words, orders, and severities made from known primitives; and
  \item new primitive types supplied only through their own single-shift
  fingerprint and behavior points, with no training pair containing that
  primitive.
\end{enumerate}
The second is harder: averaged across the $S$, $P$, and $D$ primitive-holdout
tests, the five-parameter $L_p$ model slightly outperformed \MB{} in BPE
($0.0272$ versus $0.0294$), although \MB{} was best on held-out $S$.  This is
important evidence that the simple saturation law can be preferable when
interaction data are extremely scarce.

% =============================================================================
\section{What is novel, stated precisely}
% =============================================================================

None of the ingredients is individually new mathematics: $L_p$ norms,
semigroups, upper-triangular matrices, Heisenberg groups, commutators,
attention, and neural encoders are established tools.  The project's novel
engineering/scientific contribution is their particular combination for
\emph{behavior forecasting of a frozen classifier under ordered distribution
shifts}:
\begin{enumerate}
  \item behavior deltas are accumulated through a learned signed $L_p$ law to
  encode multi-shift saturation;
  \item symmetric and antisymmetric pair terms are algebraically separated in
  parallel $3\times3$ blocks;
  \item the antisymmetric term is parameterized by positive amplitudes and
  potential differences, enabling transitive order prediction;
  \item composition is an ordered matrix product, so order sensitivity and
  associativity coexist by construction; and
  \item the output is a five-metric behavior forecast without evaluating the
  held-out composed images at prediction time.
\end{enumerate}
The strongest evidence for the combination comes from component ablations:
$p\equiv1$ damages overall prediction, while setting the commutator to zero
damages reverse-order prediction.  The novelty claim should be understood as a
claim about this modeling synthesis and experimental task, not as the invention
of the underlying algebraic objects.

% =============================================================================
\section{Assumptions, limits, and falsifiable boundaries}
% =============================================================================

\subsection{The problem is impossible without assumptions}

Primitive marginals alone cannot uniquely determine arbitrary compositions.
One can construct two worlds with identical $(v_X,\dB_X)$ and
$(v_Y,\dB_Y)$ but completely different $\dB_{X\to Y}$.  Pair calibration and
the algebraic inductive bias rule out some of these worlds; they do not prove a
universal physical law for every possible shift.

\subsection{Pairwise sufficiency}

The Heisenberg state summarizes all-pairs interactions.  A process with an
irreducible three-way term that is invisible in every primitive and pair can
still defeat pair-only training.  Regime III directly tests whether a small
amount of triple supervision helps; it does in the reported benchmark.

\subsection{Fixed classifier and reference distribution}

Predictions concern a fixed classifier, preprocessing pipeline, clean reference
distribution, and behavior definition.  Updating the classifier or changing
the clean population can change the feature fingerprint and learned algebra.

\subsection{Finite-dimensional summary loss}

The 518-vector stores first and second moments, not the complete shifted image
distribution.  Distinct distributions can share these moments.  The encoder
cannot distinguish them unless their primitive behavior vector or other
included moments differ.

\subsection{Associativity scope}

The upper-triangular operator product is exactly associative.  The neural
binary readout need not be, and the three-shift implementation explicitly
averages two decoding trees.  Claims of ``zero associativity error'' refer to
the matrix core diagnostic unless a statement explicitly includes the
symmetrized triple definition.

\subsection{Empirical scope}

The experiments support the method on the tested synthetic suites, corruption
families, architectures, and datasets.  They do not establish a theorem that
all real-world distribution shifts obey the learned $L_p$--Heisenberg law.
The counterfactual $H$ experiment strengthens the mechanism claim by separating
high-frequency structure from raw pixel variance, but it remains one designed
counterfactual.

% =============================================================================
\section{A concise end-to-end algorithm}
% =============================================================================

\textbf{Calibration.}
\begin{enumerate}
  \item Freeze a classifier and choose a clean calibration sample.
  \item For each primitive/severity, transform the images, run the frozen
  classifier, and compute $v_X$ by Equation~\eqref{eq:518} and $\dB_X$ by
  Equation~\eqref{eq:behavior}.
  \item Measure selected pair conditions for training only.
  \item Fit Stage-0 exponents on observed pairs, then optimize
  Equation~\eqref{eq:loss}.
\end{enumerate}

\textbf{Prediction for an unseen word.}
\begin{enumerate}
  \item Retrieve only $(v_{X_i},\dB_{X_i})$ for each primitive in order.
  \item Encode $z_i=E(v_{X_i})$.
  \item Build every primitive's 260 parallel $3\times3$ blocks.
  \item Left-multiply them as $T_m\cdots T_1$.
  \item Extract $b_{L_p}$, $\bar z$, $\Gamma_z$, and $\Omega_{\mathrm{ord}}$.
  \item Apply Equation~\eqref{eq:mbreadout}; for a triple, apply the exact
  symmetrized readout in Equation~\eqref{eq:jordan}.
  \item If fingerprints are out of support, apply the support safeguard.
\end{enumerate}

No image from the unseen composed condition is needed in these prediction
steps.

% =============================================================================
\section{Code-to-math map and reproducibility}
% =============================================================================

\begin{longtable}{>{\raggedright\arraybackslash}p{0.35\textwidth}
                  >{\raggedright\arraybackslash}p{0.57\textwidth}}
\toprule
Source file/function & Mathematics in this tutorial\\
\midrule
\endfirsthead
\toprule
Source file/function & Mathematics in this tutorial\\
\midrule
\endhead
\path{src/corruptions_v2.py} & Equation~\eqref{eq:imageword} and the six image transformations.\\
\path{encoder_stats_from_feats} in \path{src/experiment_v2.py} & Exact construction of Equation~\eqref{eq:518}.\\
\path{behavior_v2_from} & Margin, drift, ECE, per-class accuracy, accuracy, and confidence.\\
\path{EncoderV2}, \path{DecoderV2} & The shared $518\to64$ encoder and $64\to5$ decoder.\\
\path{OperatorProductComp.T} & Primitive translation, symmetric, and directional blocks.\\
\path{OperatorProductComp._apply_s0} & Extraction of $b_{L_p}$, $\bar z$, $\Gamma_z$, and $\Omega$.\\
\path{OperatorProductComp._readout} & Equation~\eqref{eq:mbreadout}.\\
\path{OperatorProductComp.seq_forward} & Ordered product, triple tree symmetrization, and Equation~\eqref{eq:jordan}.\\
\path{train_method} & Stage-0 $p$ fitting and Equation~\eqref{eq:loss}.\\
\path{predict} & Primitive-only inference and out-of-support safeguard.\\
\path{src/metrics_v2.py} & BPE, AccMAE, rank, and order metrics.\\
\path{src/reviewer_suite_v3.py} & V3 ablations, budget curves, primitive holdouts, external datasets, counterfactual $H$, and hierarchical statistics.\\
\bottomrule
\end{longtable}

To verify committed artifacts:
\begin{lstlisting}[language=bash]
python3 reproduce.py --verify
\end{lstlisting}
To rerun the synthetic suite or reviewer suite:
\begin{lstlisting}[language=bash]
python3 reproduce.py --synthetic
python3 reproduce.py --v3
\end{lstlisting}
Compile this tutorial, on a machine with a standard LaTeX distribution, with:
\begin{lstlisting}[language=bash]
pdflatex MATHEMATICAL_TUTORIAL.tex
pdflatex MATHEMATICAL_TUTORIAL.tex
\end{lstlisting}
The second pass resolves the table of contents and cross-references.

% =============================================================================
\section{Final mental model}
% =============================================================================

The project can be remembered as four layers:
\begin{enumerate}
  \item \textbf{Measure each primitive:} turn shifted images into a
  518-dimensional unlabeled fingerprint and a five-dimensional measured
  behavior delta.
  \item \textbf{Saturate magnitudes:} map behavior deltas through a signed
  power transform, add them, and invert the transform.
  \item \textbf{Separate interaction from order:} use one Heisenberg block for
  symmetric pair products and another for an antisymmetric potential
  difference.
  \item \textbf{Compose by multiplication:} multiply primitive operators in
  physical order and decode the resulting state into five forecast numbers.
\end{enumerate}
The $L_p$ component answers ``how large should the combined degradation be?'';
the symmetric channel answers ``how do these shifts overlap?''; and the
antisymmetric channel answers ``which one came last, and did it erase or
preserve what came before?''

\begin{thebibliography}{9}
\bibitem{repo-v2}
\textit{Shift Algebra v2: Revised Research Plan}, repository file
\path{PLAN_V2.md}.

\bibitem{repo-results}
\textit{Shift Algebra v2 Results}, repository file
\path{RESULTS_V2.md}.

\bibitem{repo-v3}
\textit{Shift Algebra V3: Comprehensive Reviewer-Proofing Suite}, repository
file \path{REVIEWER_PROOFING_V3.md}.

\bibitem{heisenberg}
The matrix law used here is the standard real three-dimensional Heisenberg
group law, specialized coordinatewise and augmented with translation
accumulators.
\end{thebibliography}

\end{document}
